\section{Geometry of the forestry state}\label{geometry-of-the-forestry-state}

\subsection{1. The marked tree space}\label{the-marked-tree-space}

Let \texttt{(D,d\_D)} be a compact geographic metric space; diameter is denoted by \texttt{d}.
A model instance may take \texttt{D} to be a region,
country, estate, or bounded study landscape. A particular forest is represented
by where the standing-tree measure has support inside \texttt{D}.

Let:

\begin{itemize}
\tightlist
\item
  \texttt{G} be a finite species/genetic-class set;
\item
  \texttt{Q\_T} be a compact tree-quality space;
\item
  \texttt{Q\_L} be a compact log-quality space;
\item
  \texttt{M} be a compact management-state space.
\end{itemize}

To preserve positivity of diameter and height, the paper writes size in
log-coordinates $s=(s_d,s_h)=(\log d,\log h)$. A tree state is

\[
\xi=(\ell,a,s,g,q,m)\in
E_T=D\times\mathbb R_+\times\mathbb R^2\times G\times Q_T\times M.
\]

Age remains an explicit coordinate. Two trees can have similar diameter and height but different age, growth
potential, mortality risk, or management timing. If a future reduced model
shows that age can safely be omitted, that is an empirical model reduction.

\subsection{2. The log, processor, and operator spaces}\label{the-log-processor-and-operator-spaces}

A harvested log is represented as

\[
\lambda=(\ell,d_s,d_l,L_L,g,q_L,w)\in E_L,
\]

where \texttt{d\_s} and \texttt{d\_l} are small- and large-end diameters, \texttt{L\_L} is length, \texttt{q\_L}
is quality, and \texttt{w} is moisture or another density-relevant state.

Log mass is

\[
m_L(\lambda)=\varrho(g,q_L,w)V_L(d_s,d_l,L_L,q_L).
\]

Density and recoverable volume depend on state, so the model uses the mass map
above in place of a universal cubic-metre-to-tonne constant.

A processor state records location, normalized desired procurement flow,
specification parameters, utilization, and a sourcing-radius parameter. An
operator state records location, capacity, productivity, equipment/fuel
intensity, cost, and compliance/quality information.

These state spaces are assumed Polish. That gives the measurable and topological
structure needed for probability measures, weak convergence, selection theorems,
and dynamic programming.

\subsection{3. Why discrete biological marks are stratified}\label{why-discrete-biological-marks-are-stratified}

Species and genetic class are discrete marks. If
\texttt{g=1} and \texttt{g=2} are species labels, interpolation through \texttt{g=1.5} has no
biological meaning.

The paper therefore identifies the tree space with a finite disjoint union of
continuous strata:

\[
E_T=\bigsqcup_{g\in G}E_T^g.
\]

A standing-tree measure can be written as a vector

\[
\mu=(\mu^g)_{g\in G}.
\]

Continuous growth occurs within a biological stratum. A species change or
reclassification, if scientifically meaningful, must be represented by an
explicit discrete kernel. A continuous species axis has no biological meaning.

\subsection{4. Finite measure versus normalized probability law}\label{finite-measure-versus-normalized-probability-law}

The persistent tree state is

\[
\mu_t\in\mathcal M_+(E_T).
\]

Its total mass \texttt{mu\_t(E\_T)} matters. Harvest and mortality remove mass;
recruitment adds mass.

When total mass is positive, define the normalized law

\[
\bar\mu_t=\frac{\mu_t}{\mu_t(E_T)}.
\]

These two objects answer different questions:

\begin{itemize}
\tightlist
\item
  \texttt{mu\_t} says both \textbf{how much} tree mass/count is present and \textbf{how it is
  distributed};
\item
  \texttt{bar\_mu\_t} says \textbf{what the composition looks like conditional on being in the
  population}.
\end{itemize}

Normalizing too early destroys information. A plantation with 100 surviving
trees and one with 1,000 surviving trees can have the same normalized size
composition.

\subsection{5. Wasserstein geometry}\label{wasserstein-geometry}

For probability measures with finite second moment, the quadratic Wasserstein
distance is

\[
W_2^2(\mu,\nu)=\inf_{\gamma\in\Pi(\mu,\nu)}
\int d(x,y)^2\,\gamma(dx,dy).
\]

It asks for the least quadratic cost of transporting one probability law into
another. This makes it natural for growth in continuous state coordinates: mass
moves from smaller to larger sizes, or through quality space, while the total
probability remains one.

The formal tangent geometry gives the continuity equation

\[
\partial_t\mu_t+\nabla\cdot(v_t\mu_t)=0.
\]

That equation conserves total mass. The full forestry state also requires mass
change.

\subsection{6. Hellinger-Kantorovich / unbalanced transport geometry}\label{hellinger-kantorovich-unbalanced-transport-geometry}

A mass-changing continuity equation has the form

\[
\partial_t\mu_t+\nabla\cdot(v_t\mu_t)
=\alpha_t\mu_t+B_t,
\]

where \texttt{alpha\_t} is a signed reaction rate and \texttt{B\_t} is an external source
measure. Mortality or harvest corresponds to negative reaction; planting or
recruitment is a source.

Hellinger-Kantorovich geometry is designed for exactly this combination of
transport and reaction. It can compare finite measures with different total
mass.

The entropy-transport and dynamic formulations used here follow
\citet{liero2018entropy} and \citet{chizat2018dynamic}. Entropic approximations
can be computed by the scaling algorithms of \citet{chizat2016scaling}.

The conceptual split is therefore:

\begin{Shaded}
\begin{Highlighting}[]
\NormalTok{Wasserstein:     composition of a normalized surviving population}
\NormalTok{HK/unbalanced:   physical state when mass can appear or disappear}
\end{Highlighting}
\end{Shaded}

Each geometry therefore serves a separate state representation.

\subsection{7. A useful normalized-selection identity}\label{a-useful-normalized-selection-identity}

Diffusion and selection produce different terms. Consider a simplified finite
measure satisfying

\[
\frac{d}{dt}\langle\phi,\mu_t\rangle
=\langle L\phi,\mu_t\rangle-\langle\lambda\phi,\mu_t\rangle,
\]

with state-dependent removal hazard \texttt{lambda}. Let \texttt{M\_t=mu\_t(E)} and
\texttt{bar\_mu\_t=mu\_t/M\_t}. Differentiating the normalized expectation gives

\[
\frac{d}{dt}\langle\phi,\bar\mu_t\rangle
=\langle L\phi,\bar\mu_t\rangle
-\operatorname{Cov}_{\bar\mu_t}(\phi,\lambda).
\]

The covariance term is a selection effect: if large trees are harvested at a
higher rate, the composition of survivors changes even without any biological
shrinkage. Recruitment adds further composition terms.

The representative-survivor law determines composition; the finite measure also
determines population mass.

\subsection{8. Linear convexity versus geodesic convexity}\label{linear-convexity-versus-geodesic-convexity}

\texttt{M\_+(E)} is a convex cone under ordinary mixtures

\[
(1-\theta)\mu_0+\theta\mu_1.
\]

This is \textbf{linear convexity}. It is useful for resource measures and relaxed
controls.

Wasserstein or HK spaces may also have geodesics. A functional can be convex
along those geodesics, which is a different statement. Geodesic convexity is
relevant to transport and gradient-flow methods; linear convexity is relevant to
ordinary convex programs.

Hard grade thresholds, species labels, regime switches, and impulses preclude a
global geodesic-convexity assumption.

\subsection{9. Finite particle approximation and tree identity}\label{finite-particle-approximation-and-tree-identity}

A finite interacting system can be written

\[
X_t^{i,N},\quad i=1,\ldots,N,
\]

with empirical law

\[
\mu_t^N=\frac1N\sum_{i=1}^N\delta_{X_t^{i,N}}.
\]

Here:

\begin{itemize}
\tightlist
\item
  \texttt{N} is the finite approximation/population size;
\item
  \texttt{i} distinguishes particles in that finite representation;
\item
  the limiting law \texttt{mu\_t} is unlabelled.
\end{itemize}

This has an important implementation consequence. Repeated-tree field data may
use a stable identifier such as \texttt{(plot,\ tree\_id)}, but that identifier can remain
in a data registry. It only belongs inside the mathematical state if future
dynamics or payoffs depend directly on identity. Observable and latent attributes
remain the coordinates of the limiting state.

\subsection{10. Control topology and relaxed controls}\label{control-topology-and-relaxed-controls}

For a compact action space \texttt{A}, strict controls are progressively measurable
\texttt{A}-valued processes. A useful metric is

\[
\Delta(\alpha,\beta)=\mathbb E\int_0^T d_A(\alpha_t,\beta_t)dt.
\]

If the strict action set is non-convex, a relaxed control replaces an action by
a probability measure on actions, an element of \texttt{P(A)}. Because \texttt{A} is compact,
\texttt{P(A)} is weakly compact and convex.

Relaxation supplies compactness and convexity for existence arguments. Physical
categories remain discrete whenever mixtures lack a physical interpretation.

\subsection{11. Implementation interpretation}\label{implementation-interpretation}

An implementation should preserve the following structure:

\begin{itemize}
\tightlist
\item
  keep mass separate from normalized composition;
\item
  keep discrete biological labels discrete;
\item
  preserve the variables needed for processor compatibility;
\item
  keep observational confidence outside physical state unless it is itself a
  controlled physical variable;
\item
  maintain stable tree IDs in a data layer when repeated measurements require
  them.
\end{itemize}

The next chapter places stochastic dynamics on these spaces and constructs the
map from living trees to harvested logs.
